% !TeX root = ../main.tex
% 修改摘要内容。
\begin{abstract}
自主量子化学计算需要将科学需求、计算执行与结论验证连接为一致的决策过程。本文以 BG6022 的三个版本为研究对象，依次讨论可追溯固定工作流（V2）、自然语言驱动的工具组合（V3）和证据约束的代码生成（V4）三种思路，并重点研究 V4 的方法设计与实验表现。V2 通过水和乙醇的优化、频率与单点计算建立可复核的计算链；V3 将评价扩展至有限科学任务和对话情境，其固定科学矩阵通过 11/11 个单元，核心基准通过 18/18 个预定义情境。上述成果分别受人工批准、任务定义与对话脚本约束。V4 进一步将任务特定算法交由生成的 Python 和 ORCA 输入表达，以单决策循环、目标与计划分离、统一执行入口及累计预算约束开放任务。早期乙醇案例实现自主优化与频率分析，含人工元数据协助的正己烷案例获得四个候选的单点能量排序。按 ORCA 正常结束判定计算成功，V4 已在乙醇优化与频率、正己烷候选能量、正丁烷部分扫描点、丙酮激发态、气相与水环境能量比较以及 HCN 屏蔽张量六类任务中取得数值结果；甲醛极化率尚无结果。原七任务冻结批次的 8 次 ORCA 启动中，7 次计算成功（87.5\%）；后续成功案例按各自批次和实际理论水平分别报告。乙醇优化获得 21 个正频率且无虚频，两种基组下水环境相对气相的电子能差分别为 $-18.276507$ 和 $-16.607972$ kJ mol$^{-1}$。研究表明，生成代码能够支持多种量子化学任务的真实执行，几何对应、结果提取与完整交付仍需进一步完善。三个版本反映研究问题的递进，现有数据不构成共同基准上的性能排序。
\end{abstract}
